%!TEX root = newmain.tex

To set the stage for our translation-based verification method, we
briefly review the relevant background. Section \ref{subsect:stipula}
presents the {\Stipula} language, emphasising its constructs for
assets, states, and timed events.  Section \ref{subsect:key}
introduces the {\KeY} system, which we employ as the verification
backend.

\subsection{{\Stipula}}
\label{subsect:stipula}

A {\Stipula} contract consists of a set of parties, states, assets,
fields, and a set of functions and events. The declaration of a
contract is defined in Figure~\ref{tab:stipula-syntax}, where
\stipula{C} is the name of the contract, $\vect{\tt h}$ and
$\vect{\tt x}$ are the \emph{assets} and \emph{fields}, respectively,
$\vect{\tt A}$ are the \emph{parties}.
%
\begin{figure}[t]
  {\small
    \begin{tabular}{l@{\quad}l}
      $
      \begin{array}{l}
        \stipula{stipula C \{}
        \\
        \qquad \stipula{assets} \; \vect{\tt h}
        \\
        \qquad \stipula{fields} \; \vect{\tt x} 
        \\
        \qquad \stipula{agreement(} \vect{\tt A} \stipula{) \{  }
        \\
        \qquad \qquad \vect{{\tt A}_1} \; : \; \vect{{\tt x}_1}
        \\
        \qquad \qquad \cdots 
        % \quad {\tt //} \; \bigcup_{i \in 1..n} \vect{{\tt A}_i} \subseteq \vect{{\tt A}}   \; , \quad  \bigcup_{i \in 1..n} \vect{{\tt x}_i} \subseteq\vect{{\tt x}} \; , \quad	\bigcap_{i \in 1..n} \vect{{\tt x}_i} = \varnothing
        \\
        \qquad \qquad  \vect{{\tt A}_n} \; : \; \vect{{\tt x}_n}
        \\
        \qquad \stipula{\}} \tostate \Qwithat  
        \\
        \qquad     F
        \\
        \stipula{\}}
      \end{array}
      $
      &
        $
        \begin{array}{l@{\quad}l@{\quad}l}
          \mathit{Functions} & F \; ::= & \zero  \quad | \quad   \Qwithat \,\A : \, {\f}(\vect{\tt y})[\vect{\tt k}]\,(E)\{
                                          \,S\, ; \, W\,\} \,\tostate\, \Qwithat'\,  F
          \\
          \mathit{Prefixes} & P \; ::= & E \, \send \, {\x} \quad | \quad E \send \, \A \quad | \quad E \lolli {\h}, {\h}' 
\\
                             & & |\quad E \lolli {\h} ,  \A 
          \\
          \mathit{Statements} & S \; ::= &  \zero  \quad | \quad P \; S \quad | \quad \ifte{E}{S}{S} \, S
          \\
          \mathit{Events} &  W \; ::= & \zero  \quad | \quad \stipula{now + t} \,\event\, \Qwithat\, \{ \, S \,\} \,\tostate\, \Qwithat'\,  W 
          \\
          \mathit{Expressions }& E \; ::= &  v \quad | \quad \stipula{now} \quad | \quad X    
                                            \quad | \quad   E \, \stipula{op} \, E \quad | \quad  \stipula{nop} \, E
%\\
%\mathit{Values} & v ::= &  n \quad |\quad \stipula{false} \quad |\quad \stipula{true} \quad | \quad  s %\quad | \quad \Time \quad |\quad a 
          \\
          \\
          \\
          \\
\end{array}
        $
    \end{tabular}
    \vspace{-.6cm}
  }
  \caption{\label{tab:stipula-syntax} Syntax of {\Stipula}}
  \vspace{-.5cm}
\end{figure}
%
The \stipula{agreement} construct declares the parties that set the
initial value of the fields. For example, if the agreement has three
parties $\A_1, \A_2, \A_3$ and the contract has two fields
$\x_1, \x_2$, if it declares $\A_1: \x_1$ and $\A_2, \A_3 : \x_2$ then
$\x_1$ will be set by $\A_1$ and $\x_2$ will be set upon agreement on
the value between $\A_2$ and $\A_3$.
%
When the agreement is concluded, the initial state of the contract is
set and the parties may invoke a function in $F$.

A \emph{function} $\f$ declared as
$\Qwithat\,\A :\,{\f}(\vect{\tt y})[\vect{\tt k}]\,(E)\{\,S;\,W\,\}\,\tostate\,\Qwithat'$ 
can be invoked by a party $\A$ if the contract is in state $\Q$ and the guard $E$ is
\emph{true}.  The names $\vect{\tt y}$ and $\vect{\tt k}$ are the
formal parameters of {\f}; they are kept separate because
$\vect{\tt y}$ are \emph{field} values while $\vect{\tt k}$ are
\emph{asset} quantities.

Function bodies are \emph{statements} $S$ followed by \emph{events}
$W$.  The former include value transfers, asset movements, conditional
logic, and field assignments. {\Stipula} distinguishes between
different transfer operations: field updates and messages use the
symbol $\rightarrow$ with the usual semantics of assignment, while
asset transfers, called \emph{moves}, use the linear operator $\lolli$
to emphasise the conservation semantics. For instance, an expression
like ``$1 \lolli \stipula{wallet}, \stipula{Seller}$'' denotes
exclusive transfer of a unit in {\tt wallet} to the {\tt Seller} and,
\emph{at the same time}, the decrease by 1 of {\tt wallet}.  In
contrast ``\stipula{wallet} $\rightarrow$ \stipula{Seller}'' specifies
that the value currently stored in \stipula{wallet} is communicated to
\stipula{Seller}, while leaving the state of \stipula{wallet}
unchanged.  It is common practice in {\Stipula} to use the
abbreviations $\h\lolli\h'$ and $\h\lolli\A$, which stand for
$\h\lolli\h,\h'$ and $\h\lolli\h,\A$, respectively.
% These abbreviations improve readability and conciseness by omitting
% redundant information while preserving the underlying semantic
% structure of the corresponding move operations.

Events
$\stipula{now + t} \,\event\, \Qwithat\, \{ \, S \,\} \,\tostate\,
\Qwithat'$ define a statement
$S$ to be executed if, \emph{after} \stipula{t} time units from the
current execution, the contract is in the state $\Q$. The expression \stipula{t}
is either a non-negative number or a variable (a formal parameter or a 
field) or a variable plus a non-negative number.
If the event is
executed, the contract will transit to the state $\Q'$. The expression
\stipula{now + t} is called \emph{time guard}.  Functions and events
are generically called \emph{clauses}.

Expressions $E$ include integer constant values $v$, the keyword
\stipula{now} (that stands for a non-negative integer value), names of
assets, fields and parameters generically ranged over by $X$, standard
arithmetic and Boolean operations \stipula{op} (binary) and
\stipula{nop} (unary). Booleans are encoded by integers in the usual
way. Actually, {\Stipula} has a richer set of values (reals, strings, time and
asset values). For the time being, we stick to integer arithmetic,
where deductive verification systems exhibit a high degree of
automation. Other datatypes require corresponding support in automated
theorem provers and SMT solvers: experimentation is needed to
determine what is possible. We also want to keep the presentation simple.

% We refer to article~\cite{CrafaLSV23} for background on the design
% of {\Stipula}.
%
The semantics of {\Stipula} is defined operationally by means of a
transition relation between \emph{configurations}
$\contract {\comma} \Time$, where $\Time$ is a positive integer value
representing the system's global clock and $\contract$ is a tuple
$\C(\State\semi\ell\semi \Sigma \semi\SemEvent)$, where:
%
\begin{itemize}
\item $\C$ is the {\Stipula} contract name;
\item $\State$ is the contract's current state: either $\zero$ (for no
  state) or a contract state \state;
\item $\ell$ maps names (parties, fields, assets, function parameters)
  to values;
\item $\Sigma$ is the (possibly empty) residual of a function body or
  of an event handler, \emph{i.e.}~$\Sigma$ is either $\zero$ or a
  term $S\ W \tostate \texttt{@}\state$;
\item $\SemEvent$ is a (possibly empty) multiset of \emph{pending
    events} to be executed. We let $\SemEvent$ be $\zero$ when there
  are no pending events, or
  %
  \begin{equation}
    \label{eq:psi}
    {\Time}_1 \,\event_{\ev_1}\, \Qwithat_1\, \{ \, S_1 \,\}
    \,\tostate\, \Qwithat_1' \quad|~ \cdots ~|\quad {\Time}_n \,\event_{\ev_n}\,
    \Qwithat_n\, \{ \, S_n \,\} \,\tostate\, \Qwithat_n'\enspace,
  \end{equation}
  %
  where $\Time_i$ are the values of the expressions \stipula{now +}
  \stipula{t}$_i$ of the corresponding events.  The indices $\ev_i$
  are introduced by rule~\rulename{function} when the body of a
  function is incorporated into the configuration and denotes a unique
  identifier
  % . Each such index denotes the unique integer identifier
  statically assigned to the event.
\end{itemize}

The transition relation of {\Stipula} is
$\contract {\comma} \Time \lred{\mu} \contract' {\comma} \Time'$,
where $\mu$ may be empty, or an initial agreement, or
$A:\f(\wt{u})[\wt{v}]$ (in case of function invocations) or an event.
%
Its formal definition is given in Table~\ref{tab:controlrules} for the
control transition rules, see \cite{CrafaLSV23} for the agreement and
statement rules. The former rule set requires a specific
{\JML}/{\Java} encoding, whereas the latter are translated in a
standard way and are inessential for the understanding of this paper.
%
\begin{table}[t]
{\small
$
\begin{array}{c}
\mathrule{Function}{
	\begin{array}{c}
	\texttt{@}\state\,\A\,\texttt{:}\,\f\texttt{(} \wt{\mbox{{\tt y}}} \texttt{)[} \wt{\mbox{{\tt k}}} \texttt{]\,(} E {\tt )}\,{\tt \{}\,S\,W\,{\tt \}} \tostate \texttt{@}\state' \in \C
	\\
	\nored{\SemEvent}{\Time}
	%\quad \vect{v} \geq 0 
	\\
	\ell(\A) = A
	\quad
	\ell' = \ell [ \wt{\mbox{{\tt y}}} \mapsto \wt{u}, \wt{\mbox{{\tt k}}} \mapsto \wt{v} ]
	\quad
	\sem{E}_{\ell'} \neq 0
	\end{array}
	}{
	\C(\state \semi \ell \semi \zero \semi \SemEvent) \comma \Time 
	    \lred{A : \texttt{f}(\wt{u})[\wt{v}]}
	\C(\state \semi \ell' \semi S\,W \tostate \texttt{@}\state' 
	\semi \SemEvent) \comma \Time
}	
\quad
\mathrule{State-Change}{
    \sem{W\subst{\Time}{\texttt{now}}}_{\ell} = \SemEvent'
    }{
    \begin{array}{l}
	\C(\state \semi \ell \semi \zero \,W\tostate \texttt{@}\state' \semi \SemEvent) 
	\comma \Time 
	\\
	\qquad \lred{}
	\C(\state' \semi \ell \semi \zero \semi \SemEvent' ~|~\SemEvent) 
	\comma \Time
	\end{array}
}
\\
\\
\mathrule{Event-Match}{
	\begin{array}{c}
	\SemEvent = \Time \event_{\ev_\kappa} \texttt{@}\state~\{~S~\} \tostate \texttt{@}\state' ~|~ \SemEvent'
	\end{array}
	}{
	\begin{array}{l}
	\C(\state \semi \ell \semi \zero \semi \SemEvent) \comma \Time 
%	\\
%	\qquad 
\lred{\ev_\kappa}
	\C(\state \semi \ell \semi S \tostate \texttt{@}\state' \semi \SemEvent') 
	\comma \Time  
	\end{array} 
}
\quad
\mathrule{Tick}{
	\nored{\SemEvent}{\Time}
	}{
	\begin{array}{l}
	\C(\state \semi \ell \semi \zero \semi \SemEvent) \comma \Time 
	\\
	\qquad 
\lred{} 
	\C(\state \semi \ell \semi \zero \semi \SemEvent\downarrow_{\Time}) \comma \Time + 1
	\end{array}
	}
\\
\\
\end{array}
$
}
\caption{\label{tab:controlrules} The control transition rules of {\Stipula}}
\vspace{-.6cm}
\end{table}

The formal definition of
$\contract {\comma} \Time \lred{\mu} \contract' {\comma} \Time'$ uses
the following auxiliary predicates and functions:
%
\begin{itemize}
\item $\sem{E}_{\ell}$ is a \emph{partial function} returning the
  value of $E$ in memory $\ell$.
  % \item $\sem{E}_\ell^\mathsf{a}$ is the partial function returning
  %   asset values. This evaluation differs from $\sem{E}_{\ell}$
  %   because it never return negative values (assets are always
  %   nonnegative)
\item $\sem{W\{^\Time/_{{\tt now}}\}}_\ell$ is the multiset of
  scheduled events obtained from the sequence $W$ by replacing every
  {\tt now} by $\Time$ and evaluating every expression in the time
  guards.
\item $\nored{\SemEvent}{\Time}$ is a predicate that is \emph{true}
  whenever $\SemEvent$ is as in equation~\eqref{eq:psi}
  % $\SemEvent = \Time_1 \event \texttt{@}\state_1 \{ S_1 \} \tostate
  % \texttt{@}\state_1' ~|~ \cdots ~|~ \Time_n \event \texttt{@}\state_n
  % \{ S_n \} \tostate \texttt{@}\state_n'$
  and, for every $1 \leq i \leq n$, $\Time_i \neq \Time$. It is
  \emph{false} otherwise.
\item
$\SemEvent\downarrow_{\Time}$ removes from $\Psi$ every event with time value less or equal to $\Time$. 
\end{itemize}

We explain the rules in Table~\ref{tab:controlrules}.
%
Rule \rulename{Function} defines invocations: the label specifies
party $A$ performing the invocation of a function named {\tt f} with
the actual parameters. The transition may occur provided (\emph{i})
the contract is in the state \texttt{Q}, admitting invocations of {\tt
  f} from $A$, (\emph{ii}) the contract is \emph{idle},
\emph{i.e.}~the contract has no statement to execute (presence of
$\zero$ in the left configuration), (\emph{iii}) the precondition $E$
is satisfied, and no event can be triggered (premise
$\nored{\SemEvent}{\Time}$). In particular, the final constraint
expresses that events \emph{have precedence} on possible function
invocations. For example, if a payment deadline is reached and, at the
same time, the payment arrives, it will be refused in favour of the
event managing the deadline.  We use two notations for party
identifiers: {\A} and $A$. The former denotes the formal name used in
the contract specification, whereas $A$ refers to the concrete party
name instantiated and fixed when the agreement is concluded.

Rule \rulename{State-Change} says that a contract changes state when
the execution of the statement in the function body terminates. At
this stage, the sequence of events $W$ is added to the multiset of
pending events once their time expressions have been evaluated, in
particular \texttt{now} is replaced by the current value of the clock.
%
Rule \rulename{Event-Match} specifies that event handlers may run
provided there is no statement to perform and the time guard of the
event has exactly the value of $\Time$.
%
Rule \rulename{Tick} defines the elapsing of time, which may happen
when there is no statement to perform and no event can be triggered.

In this paper we use the state transition models of {\Stipula}
contracts, called the \emph{underlying automata}. The states of these
automata are those of the {\Stipula} contract. The transitions
correspond to clauses and are labelled either with the function name
(the party name is always omitted, for simplicity's sake we assume
that function names are pairwise different) or with a uniquely
numbered event. By definition, these automata are \emph{finite models};
% line number (\emph{e.g.}, \stipula{ev}$_{10}$ is the event at code
% line 10).
Figures~\ref{fig:license} and~\ref{fig:deposit-fsm} show the two
underlying automata of the contracts \stipula{Licence} and
\stipula{Deposit}.
%
These contracts are drawn from~\cite{iFM25} with minor modifications, which 
are discussed below. 

\begin{example}
  Figure~\ref{fig:license} defines the \stipula{License} contract
  regulating a licensing transaction between a \stipula{Licensor} and
  a \stipula{Licensee}, with a 30 days trial period and the
  possibility to purchase or decline the license. The field {\tt code}
  stores the information for the trial period access, the asset {\tt
    token} grants the licence.  If the \stipula{Licensee} does not
  purchase the license before the trial period expires, the contract
  terminates and the {\tt token} is automatically returned to the
  \stipula{Licensee}.
  % 
  The \stipula{agreement} in Line~\ref{line:lic:agb} allows the
  parties to agree on the license \stipula{cost}. Upon agreement, the
  first state is \stipula{@Init} where
  % 
\begin{figure}[h]
  % \centering
  \quad\quad
\begin{tabular}{@{}c@{\hspace{-.8cm}}c@{}}
{\scriptsize
\begin{lstlisting}[numbers=left,frame=none,xleftmargin=2.5em,mathescape=true,language=stipula,
escapechar=§]
stipula License {
    assets balance, token
    fields cost, code
    agreement (Licensor, Licensee)(cost) { §\label{line:lic:agb}§
        Licensor, Licensee : cost
    } $\tostate$ @Init §\label{line:lic:age}§
    @Init Licensor: offer(x)[n] {
        n $\lolli$ token
        x $\send$ code
        now + 1 $\event$ @Prop { token $\lolli$ Licensor } $\tostate$ @End §\label{line:event_init}§
    } $\tostate$ @Prop 
    @Prop Licensee: activate()[b] (b == cost) {
        b $\lolli$ balance
        code $\send$ Licensee ;
        now + 30 $\event$ @Trial {
                balance $\lolli$ Licensee
                token $\lolli$ Licensor
                } $\tostate$ @End
    } $\tostate$ @Trial
    @Trial Licensee: buy()[] {
        balance $\lolli$ Licensor
        token $\lolli$ Licensee
    } $\tostate$ @End
}
\end{lstlisting}}
&
\scalebox{0.80}{
\begin{tikzpicture}[->, >=stealth', node distance=2.5cm, every state/.style={draw, minimum size=1.2cm},initial text=]
  \node[state, initial] (init) {\stipula{Init}};
  \node[state, right of=init] (first) {\stipula{Prop}};
  \node[state, right of=first] (run) {\stipula{Trial}};
  \node[state, below of=first] (end) {\stipula{End}};

  \path (init) edge node[above] {\stipula{offer}} (first)
        (first) edge node[above] {\stipula{activate}} (run)
        (first) edge node[left] {\stipula{ev}$_{1}$} (end)
        (run) edge node[left] {\stipula{buy}} (end)
        (run.south) edge node[right] {\stipula{ev}$_{2}$} (end.east);
\end{tikzpicture}
}
\end{tabular}
%\vspace{-.3cm}
\caption{The License contract in {\Stipula} and its underlying automaton.}
\label{fig:license}
\vspace{-.3cm}
\end{figure}
%
the \stipula{Licensor} may invoke \stipula{offer}, transferring a
\stipula{token} (representing the license) into escrow and a trial
\stipula{code}.  A scheduled \stipula{event} is simultaneously
registered after 1 day to reclaim the token if the \stipula{Licensee}
fails to start the trial period.  The contract then moves to
\stipula{Prop}.
% 
If no function is called in this state, then time is advanced,
eventually triggering the event at line~\ref{line:event_init}.
% 
In \stipula{Prop}, the \stipula{Licensee} can activate the trial by
paying the \stipula{cost} (observe that \stipula{b} is an asset of the
\stipula{Licensee}), which is transferred to the contract's
\stipula{balance}: the contract acts as a notary for assets that are
not finally disposed.  Now the license code is revealed to the
\stipula{Licensee} and another \stipula{event} is scheduled to handle
the expiration of the trial period: if the license is not purchased
before 30 days, the balance is refunded, the token returned to the
\stipula{Licensor}.
%
The license purchase is done by the \stipula{buy} function, which
finalises the transaction by transferring the \stipula{balance} to the
\stipula{Licensor} and assigning the \stipula{token} permanently to
the \stipula{Licensee}. 
In the corresponding contract in~\cite{iFM25}, the start and end of the 
trial period are defined during the agreement, whereas here, for simplicity, 
they are hard-coded in the contract.
\end{example}

\begin{example}
  Figure~\ref{code:deposit-stipula} defines the \stipula{Deposit}
  contract that
%
\begin{figure}[t]
\centering
\begin{tabular}{c@{\hspace{-.6cm}}c}
{\scriptsize
\begin{lstlisting}[numbers=left,frame=none,xleftmargin=2.5em,mathescape=true,language=stipula,escapechar=§]
stipula Deposit {
    asset flour
    field cost_flour
    agreement (Client, Farm)(cost_flour) {
        Client, Farm : cost_flour
    } $\tostate$ @Start
    @Start Farm : begin()[h]{
        h $\lolli$ flour;
        now + 365 $\event$ @Run { flour $\lolli$ Farm } $\tostate$ @End
    } $\tostate$ @Run
    @Run Farm : send()[k]{
        k $\lolli$ flour
    } $\tostate$ @Run
    @Run Client : buy()[w](w/cost_flour <= flour){
        (w/cost_flour) $\lolli$ flour, Client
        w $\lolli$ Farm
    } $\tostate$ @Run
}
\end{lstlisting}}
& 
\scalebox{0.8}{
\begin{tikzpicture}[->, >=stealth', node distance=2.5cm, every state/.style={draw, minimum size=1.2cm}, initial text=]
  \node[state, initial] (start) {\stipula{Start}};
  \node[state, right of=start] (run) {\stipula{Run}};
  \node[state, right of=run] (end) {\stipula{End}};

  \path (start) edge node[above, sloped] {\stipula{begin}} (run)
        (run) edge[loop above] node[above] {\stipula{buy}} (run)
        (run) edge[loop below] node[below] {\stipula{send}} (run)
        (run) edge node[above, sloped] {\stipula{ev}$_{1}$} (end);
\end{tikzpicture}
}
\end{tabular}
%\vspace{-.3cm}
\caption{The \stipula{Deposit} contract in {\Stipula} and its underlying automaton}
\label{code:deposit-stipula}
\label{fig:deposit-fsm}
\vspace{-.5cm}
\end{figure}
%
models the interactions between a \stipula{Farm} and a
\stipula{Client}.  The \stipula{Farm} deposits flour, while the
\stipula{Client} purchases and withdraws the corresponding amount at
an agreed price.  The contractual terms are enforced over a validity
period of 365 days.  The clause \stipula{send} allows the
\stipula{Farm} to deposit flour into the contract's stock; the
operation
% \stipula{h $\send$ Client} is informs the \stipula{Client} about the
% deposited amount, while
\stipula{h $\lolli$ flour} implements the deposit in the asset
\stipula{flour}.  Actual delivery occurs only with \stipula{buy()[w]},
where \stipula{(w/cost_flour) $\lolli$ flour, Client} transfers flour
from the contract's stock to the \stipula{Client} (hence
\stipula{w/cost_flour} must not be greater than \stipula{flour}) and
\stipula{w $\lolli$ Farm} represents the payment to the
\stipula{Farm}.
% In this way, {\tt send} supplies the stock, whereas {\tt buy}
% withdraws from it under payment.
In~\cite{iFM25}, the \stipula{Deposit} contract does not allow \stipula{buy} and \stipula{send} 
to start and end in the same state, as the underlying theory does not cover this case. 
Consequently, a more convoluted contract was required. In contrast, the definition 
of \stipula{Deposit} presented here is simpler, as it relies on a more expressive 
theoretical framework.
\end{example}

\stipula{License} and \stipula{Deposit} represent paradigmatic
interaction patterns that frequently arise in legal contracts and that
are clearly reflected in their underlying automata. The behavioural
structure of \stipula{License} is \emph{acyclic}: no state is
revisited during execution.  In contrast, \stipula{Deposit} exhibits a
\emph{cyclic} pattern, as it includes functions that may be invoked
repeatedly, leading to repeated traversals of certain states.  It is
important to observe that, while the semantics of {\Stipula} admits
potentially infinite cyclic behaviour (for example, \stipula{Deposit}
allows an unbounded number of interactions in a fixed time window of
365 days), in practice only finitely many transactions can
realistically occur. In Section~\ref{sec:cyclic-behavior}, we use this
observation to obtain bounded encoding of cycles
% and simpler loop invariants, which are used to implement cyclic behaviour
in the translated contracts.

\begin{definition}[Cyclic Contract, Separate Cycle]
  \label{def:determinate}
  A {\Stipula} contract is \emph{cyclic} if its underlying automaton
  has a \emph{cycle}: a sequence
  $\Q_1 \lred{\nu_1} \Q_{2} \lred{\nu_2} \cdots \lred{\nu_{\ell-1}}
  \Q_\ell \lred{\nu_\ell} \Q_{1}$ where $\Q_1,\ldots,\Q_\ell$ are
  pairwise different and $\nu_i$ are either function names or
  events. A contract has \emph{separate cycles} if any two cycles
  either have no state in common or share their initial state only.
  % $\Q_1 \lred{\mu_1} \Q_{2} \lred{\mu_2} \cdots \lred{\mu_{\ell-1}} \Q_\ell 
  % \lred{\mu_\ell} \Q_{1}$
  % of pairwise different states $\Q_1,\ldots,\Q_n$ with transitions
  % $\Q_i \lred{\mu_i} \Q_{i+1}$ and $\Q_n \lred{\mu_n} \Q_{1}$ where
  % $\mu_i$ are either function names or events.  The contract has
  % \emph{disjoint cycles} if different cycles have no state in common.
  A contract is \emph{acyclic} if it has no cycle.
\end{definition}

The contract \stipula{License} is acyclic, while \stipula{Deposit} has
separate cycles.  In what follows, we first consider acyclic contracts
and then turn to separate cyclic ones. This progression allows us to
introduce the complexities of the translation into {\JML}/{\Java} in a gradual
manner.

\subsection{Deductive Verification with the {\KeY} System}
\label{subsect:key}

The {\KeY} system \cite{DBLP:series/lncs/10001,KeYTutorial24} is a
deductive verification framework for {\Java} that combines symbolic
execution, invariant reasoning, and method contracts\footnote{Not to
  be confused with {\Stipula} contracts. In this paper, both kinds of
  contract are featured, but it should be obvious from the context
  when we mean {\Stipula} contracts and when {\JML}/{\Java} method
  contracts.} on top of a calculus for an expressive program logic.
%
The main use case of {\KeY} is formal verification of Java programs
annotated with specifications written in {\JML}. 
It provides an interactive user environment, where one can
construct correctness proofs of {\Java} methods against their {\JML}
contracts. Method contracts typically include preconditions
(\jml{requires}), postconditions (\jml{ensures}), frame
conditions (\jml{assignable}), class invariants, and auxiliary
annotations such as loop invariants and assertions. Given a
{\JML}-annotated {\Java} class, {\KeY} translates the specifications into
logical proof obligations and attempts to discharge them using its
symbolic execution engine.
%
The details of the verification process are irrelevant for the purpose
of this paper: {\KeY} is used as a \emph{black box}, the input is a
{\JML}-annotated Java file that results from translation of a
{\Stipula} contract. In principle, \emph{any} deductive verification
system with a sufficient coverage of \JML/\Java\ can be used, for
example, Krakatoa~\cite{Krakatoa07} or OpenJML~\cite{Cok21}, although
we did not try that.  In general, {\KeY} is used in interactive or
auto-active mode, however, for all case studies discussed below, the
verification is \emph{fully automatic}.

%%% Local Variables:
%%% mode: latex
%%% TeX-master: "newmain"
%%% End:
